\section{Da limitação à agenda: o que a Fase 2 corrige}
\label{sec:debate-fase2}

\begin{confronto}
``Vocês listaram com honestidade as fraquezas (Seção~\ref{sec:debate-sintese}), mas listar não é
corrigir. Sem um plano concreto, a transparência vira álibi para a inação.''
\end{confronto}

A crítica é justa e foi antecipada. A segunda fase do projeto (milestone~M8) converte três das
fronteiras assumidas em tarefas com decisão registrada (ADR-0009, 0010 e 0011, em status
\emph{proposto}) e ordem de prioridade explícita. O ponto desta seção não é apresentar
resultados --- não há ainda ---, mas mostrar que cada limitação tem um caminho de correção
\emph{já desenhado}, com os riscos mapeados.

\subsection{Diluição do choque $\to$ agregação por máximo/percentil}
\label{sec:debate-fase2-pooling}
A média sobre 30 passos atenua um choque de um dia (Seção~\ref{sec:debate-sintese}, item~4). A
correção troca a agregação por \emph{max} ou percentil alto, recuperando o pico.
\begin{tradeoff}
Maior \emph{recall} de choques pontuais ao custo de potencial perda de \emph{precision}: o máximo
é sensível também ao ruído de um único dia. Mitiga-se com o percentil (p90) em vez do máximo puro,
e mede-se o efeito real com P/R/F1 nos quatro ativos.
\end{tradeoff}
\begin{honesto}
A mudança invalida o limiar p95 atual (calibrado sobre a média): ele \textbf{tem} de ser
recalibrado sobre o novo escore. Reportar Recall de \emph{max} com limiar de \emph{mean} seria um
erro de comparação --- e está explicitamente vedado no plano (ADR-0009).
\end{honesto}

\subsection{\emph{Holdout} único $\to$ validação \emph{walk-forward}}
\label{sec:debate-fase2-cv}
O item~1 da síntese --- validação cruzada temporal --- materializa-se com \texttt{TimeSeriesSplit}
(ADR-0010), respondendo à objeção de variância levantada na Seção~\ref{sec:debate-valid-cv}.
\begin{honesto}
A validação cruzada não melhora o detector; melhora a \emph{confiança} na escolha de
hiperparâmetros. É instrumento de credibilidade, não de desempenho --- e é assim que será
reportada, sem inflar sua contribuição.
\end{honesto}
A regra anti-vazamento (ADR-0001) é preservada \textbf{por fold}: scaler reajustado só no treino
de cada corte, janelas geradas dentro da partição.

\subsection{Univariado $\to$ entrada multivariada (OHLCV)}
\label{sec:debate-fase2-multi}
Os itens~2 e~3 da síntese (anomalias de volume, hoje invisíveis) tornam-se a transição para um
tensor multivariado (ADR-0011), começando por Close+Volume.
\begin{tradeoff}
Ganha-se um canal precursor (volume) ao custo de $5\times$ a entrada, mais capacidade e mais
risco de \emph{overfitting} --- razão pela qual depende da validação \emph{walk-forward} já
disponível e começa pelo mínimo informativo, não pelo OHLCV completo.
\end{tradeoff}

\subsection{Por que nesta ordem}
\label{sec:debate-fase2-ordem}
A prioridade não é arbitrária: \textbf{(1)} a agregação é o ganho rápido de \emph{recall} que
custa poucas linhas; \textbf{(2)} a validação \emph{walk-forward} é a régua honesta para
comparar (1) contra a \emph{baseline} e calibrar (3); \textbf{(3)} o multivariado é a aposta de
maior alcance, viável apenas depois que o sinal deixa de ser diluído e os números passam a ser
confiáveis. Corrigir na ordem inversa seria medir uma aposta cara com um instrumento que ainda
esconde o que se quer medir.
